% Section 8: Extensions
\label{sec:extensions}

\subsection{Endogenous Signal Correlation: A Prisoner's Dilemma in AI Adoption}

The main model treats $\rho$ as exogenous. This section asks: why do agents converge on a common model, and can they stop? We endogenise $\rho$ by modelling fund managers' AI adoption decisions as a strategic game. Each manager has a private incentive to adopt the dominant AI model. The aggregate effect of adoption raises effective correlation above the socially optimal level, activating the fragility channels of Propositions~\ref{prop:uniqueness-boundary}--\ref{prop:bifurcation}. The answer is a prisoner's dilemma.

\subsubsection{Environment}

A continuum of fund managers indexed by $i \in [0,1]$ each choose a benchmark weight $\rho_i \in [0,1]$, representing the fraction of their portfolio signal derived from the dominant AI model. Given the choice $\rho_i$, manager $i$'s signal is $x_i = \theta + \sqrt{\rho_i}\,\eta + \sqrt{1-\rho_i}\,\xi_i$, with the same structure as~\eqref{eq:noise-decomposition}. The aggregate correlation is $\bar{\rho} = \int \rho_i \, di$; in a symmetric equilibrium, $\rho_i = \bar{\rho}$ for all $i$.

Manager $i$'s payoff has three components. First, expected trading profit from private information, following the Grossman-Stiglitz framework of Channel 2:
\begin{equation}\label{eq:pi-alpha-adoption}
  \pi_\alpha(\rho_i, \bar{\rho}) = \frac{k_P \sqrt{1-\rho_i}}{1+k_{\text{comp}} \bar{\rho}},
\end{equation}
where $\sqrt{1-\rho_i}$ captures idiosyncratic content (trading edge from private signal component) and $1/(1+k_{\text{comp}} \bar{\rho})$ captures aggregate competition (when $\bar{\rho}$ is high, many agents trade on similar signals). Second, a tracking error cost from deviating from the consensus AI signal:
\begin{equation}\label{eq:tracking-error}
  TE_i(\rho_i, \bar{\rho}) = c_{\text{TE}} (1-\rho_i) \bar{\rho},
\end{equation}
increasing in the deviation $(1-\rho_i)$ from the AI signal and in $\bar{\rho}$ (when most managers use the AI model, the benchmark is the AI signal). Third, a systemic risk externality from the amplification loop:
\begin{equation}\label{eq:systemic-cost}
  SC(\bar{\rho}) = \frac{\alpha_{\text{SC}}^2}{(1-\alpha_{\text{SC}})^2} \cdot \frac{\kappa_{\text{sys}}}{\tau} \cdot \bar{\rho},
\end{equation}
which depends on $\bar{\rho}$ (the aggregate correlation) but not on $\rho_i$ (the individual choice). Each manager bears a share of the systemic fragility cost characterised in Proposition~\ref{prop:bifurcation} regardless of her own $\rho_i$.

\subsubsection{Nash Equilibrium and Social Optimum}

Because $SC(\bar{\rho})$ depends only on $\bar{\rho}$, which each manager takes as given, the individual problem is $\max_{\rho_i} \pi_\alpha(\rho_i, \bar{\rho}) - TE_i(\rho_i, \bar{\rho})$. The first-order condition yields the best response $\rho_i^*(\bar{\rho}) = 1 - \Gamma(\bar{\rho})^2$, where $\Gamma(\bar{\rho}) = k_P / (2 c_{\text{TE}} \bar{\rho}(1+k_{\text{comp}} \bar{\rho}))$ is strictly decreasing in $\bar{\rho}$. Hence $\rho_i^*$ is strictly increasing in $\bar{\rho}$: AI adoption is a strategic complement. When more managers adopt the AI model, the tracking error cost of deviating rises and the alpha from idiosyncratic research falls, strengthening each individual manager's incentive to adopt.

\begin{proposition}[Endogenous Signal Correlation]\label{prop:endogenous-rho}
  In the AI adoption game:

  \emph{(i)} The individual best response $\rho_i^*(\bar{\rho})$ is strictly increasing in $\bar{\rho}$.

  \emph{(ii)} There exists a unique stable symmetric Nash equilibrium $\rho^{\text{NE}} \in (0,1)$, characterised by
  \begin{equation}\label{eq:nash-eq}
    k_P = 2 c_{\text{TE}} \rho^{\text{NE}} \sqrt{1-\rho^{\text{NE}}} (1 + k_{\text{comp}} \rho^{\text{NE}}).
  \end{equation}

  \emph{(iii)} The Nash equilibrium correlation exceeds the social optimum: $\rho^{\text{NE}} > \rho^{\text{SO}}$. Under the condition $\kappa_{\text{sys}} > \kappa_{\text{sys}}^* \equiv \tau(1-\alpha_{\text{SC}})^2(\rho^{\text{NE}} - \rho^*)/(\alpha_{\text{SC}}^2 \rho^{\text{NE}})$, which ensures the systemic cost term in~\eqref{eq:systemic-cost} shifts the social optimum below the bifurcation threshold, the gap is sufficiently large that $\rho^{\text{NE}} > \rho^*$.

  \emph{(iv)} All managers are strictly worse off at the Nash equilibrium than at the cooperative outcome: $U(\rho^{\text{NE}}, \rho^{\text{NE}}) < U(\rho^{\text{SO}}, \rho^{\text{SO}})$.

  \prooflabel{Proof sketch.} Part (i): $\Gamma(\bar{\rho})$ is strictly decreasing in $\bar{\rho}$, so $\rho_i^* = 1-\Gamma^2$ is strictly increasing. Part (ii): The Nash equilibrium solves $\rho = \rho_i^*(\rho)$, which after substitution gives~\eqref{eq:nash-eq}. The upper root of this equation is the stable equilibrium by best-response dynamics.\footnote{The function $H(\rho) = 2 c_{\text{TE}} \rho \sqrt{1-\rho}(1+k_{\text{comp}}\rho)$ satisfies $H(0) = H(1) = 0$ with a unique interior maximum. For $k_P$ below the maximum, $H(\rho) = k_P$ has two roots; the upper root is stable.} Part (iii): The social planner maximises $W(\rho) = U(\rho, \rho) = \pi_\alpha(\rho, \rho) - TE(\rho, \rho) - SC(\rho)$, internalising the systemic risk externality~\eqref{eq:systemic-cost} and the competition externality. The social planner's first-order condition includes three additional negative terms absent from the individual's problem: a competition externality, a tracking error externality, and the systemic risk externality $-[\alpha_{\text{SC}}^2/(1-\alpha_{\text{SC}})^2]\kappa_{\text{sys}}/\tau$. At $\rho^{\text{NE}}$, the social marginal cost exceeds the social marginal benefit, so $\rho^{\text{SO}} < \rho^{\text{NE}}$. Part (iv): Since $\rho^{\text{SO}}$ maximises $W(\rho)$ and $\rho^{\text{NE}} \neq \rho^{\text{SO}}$, strict concavity of $W$ implies $W(\rho^{\text{SO}}) > W(\rho^{\text{NE}})$. $\square$
\end{proposition}

Proposition~\ref{prop:endogenous-rho}(iv) establishes the prisoner's dilemma. Each fund manager individually benefits from adopting the dominant AI model: it reduces tracking error and provides a signal that, while correlated, is individually precise. Aggregate adoption raises systemic fragility through all three channels. The tracking error motive creates a coordination trap. Once most managers use the AI model, deviating to proprietary research is individually costly, even though collective deviation would improve stability. No manager internalises the systemic externality. Every manager defects.

The Nash equilibrium correlation is increasing in $c_{\text{TE}}$ (stronger career concerns) and decreasing in $k_P$ (better private alpha). The wedge $\rho^{\text{NE}} - \rho^{\text{SO}}$ is strictly increasing in $\alpha_{\text{SC}}$: stronger coordination externalities widen the gap between private and social incentives. These conditions describe the current trajectory of the financial sector: increasing cost of fundamental research, strong peer-performance pressure, and convergence on a small number of dominant AI platforms.

\subsection{The Diversity Mandate}

A regulator observing the prisoner's dilemma characterised in Proposition~\ref{prop:endogenous-rho} can implement a portfolio diversity constraint $\rho_i \leq \rho_{\max}$ for all $i$.

\begin{proposition}[Optimal Diversity Mandate]\label{prop:diversity-mandate}
  In the presence of the prisoner's dilemma in Proposition~\ref{prop:endogenous-rho}, a diversity mandate $\rho_i \leq \rho_{\max}$ achieves:

  \emph{(i) Welfare improvement:} For any $\rho_{\max} \in (\rho^{\text{SO}}, \rho^{\text{NE}})$, the mandate strictly improves welfare relative to the Nash equilibrium: $W(\rho_{\max}) > W(\rho^{\text{NE}})$.

  \emph{(ii) First-best achievability:} Setting $\rho_{\max} = \rho^{\text{SO}}$ achieves the cooperative welfare level $W(\rho^{\text{SO}})$, provided the individual rationality constraint $U(\rho^{\text{SO}}, \rho^{\text{SO}}) \geq U_0$ is satisfied.

  \emph{(iii) Systemic safety:} Setting $\rho_{\max} < \rho^*$ (the bifurcation threshold of Proposition~\ref{prop:bifurcation}) ensures the integrated three-channel system remains below its fragility threshold.

  \prooflabel{Proof sketch.} Part (i): By Proposition~\ref{prop:endogenous-rho}(iii), $\rho^{\text{NE}} > \rho^{\text{SO}}$, so the mandate binds for any $\rho_{\max} < \rho^{\text{NE}}$. Since $W(\rho)$ is increasing for $\rho < \rho^{\text{SO}}$ and decreasing for $\rho > \rho^{\text{SO}}$, any $\rho_{\max}$ closer to $\rho^{\text{SO}}$ than $\rho^{\text{NE}}$ improves welfare. Part (ii): At $\rho_{\max} = \rho^{\text{SO}}$, the mandate implements the cooperative outcome that maximises $W$. Part (iii): From Proposition~\ref{prop:bifurcation}, $\rho > \rho^*$ implies the integrated system is in the fragile region; bounding $\rho_{\max} < \rho^*$ ensures stability. $\square$
\end{proposition}

The diversity mandate can be implemented through model diversity requirements (minimum number of distinct AI models or data sources), correlation-based capital charges, or algorithmic transparency requirements. The practical challenge is observability: direct enforcement requires observing each manager's AI model weight $\rho_i$ through portfolio disclosure or model audit; if only aggregate $\bar{\rho}$ is observable, individual-level reporting or random audits may be required.

\subsection{Discussion of Dynamic Extensions}

The main model is static. Two natural dynamic extensions merit future investigation. First, a model of evolving $\rho$ over time would characterise whether the bifurcation threshold $\rho^*$ is approached gradually (providing a regulatory warning window) or crossed discontinuously. Second, a model of architecture heterogeneity would endogenise the distribution of AI models in use, extending the prisoner's dilemma to a richer game in which different architectures have partially independent errors. Both extensions share a common feature: $\rho^*$ provides a natural regulatory target, and the diversity mandate of Proposition~\ref{prop:diversity-mandate} provides the instrument.

The theoretical analysis is now complete. We turn in Section~\ref{sec:empirics} to motivating evidence before drawing policy conclusions in Section~\ref{sec:conclusion}.
% --- Clarity pass complete: 2026-03-10 ---
% --- Flow pass complete: 2026-03-11 ---
% --- Technical audit complete: 2026-03-11 ---
